
Si \(n\ge m\),
alors \(\beta_n\) est surjectif par 1.3 (ii); donc \(\alpha_{n+1}\) est surjectif. Par conséquent, \(H^1(F(n-1))\xrightarrow{\sim}H^1(F(n))\xrightarrow{\sim}\ldots\xrightarrow{\sim}0\). D'autre part, si \(\alpha_n\) n'est pas surjectif, alors \(h^1(F(n-1))>h^1(F(n))\); donc \(n\mapsto h^1(F(n))\) est \uline{strictement} décroissant depuis \(h^1(F(m-1))\), jusqu'au moment de devenir \(0\); d'où (iii).

\textbf{Définition 1.5.} \emph{Soient \(F\) un faisceau cohérent sur \(X\), \(r\) un entier \(\ge\operatorname{dim{.}supp}(F)\) et \((b)=(b_0,\ldots,b_r)\) une suite de \(r+1\) entiers. Alors \(F\) est dit un \emph{\((b)\)-faisceau} s'il satisfait à un des deux ensembles équivalents suivants de conditions}~:

\((1.5.1)\)
\begin{enumerate}[label=(\roman*)]
\item \emph{En chaque point du support de \(F\), \({O}_X(1)\) est engendré par \(H^0({O}_X(1))\).}
\item \emph{\(h^0(F(-1))\le b_0\).}
\item \emph{(Lorsque \(r\ge1\)), il existe une section \(\sigma\in H^0({O}_X(1))\) qui définit une suite exacte \(0\to F(-1)\to F\to G\to0\) telle que \(G\) soit un \((b_1,\ldots,b_r)\)-faisceau.}
\end{enumerate}

\((1.5.2)\) \emph{Il existe une suite \(F\)-régulière de \(r\) sections \(\sigma_1,\ldots,\sigma_r\in H^0({O}_X(1))\) telle que, pour \(i=0,\ldots,r\), si \(F_i\) désigne la restriction de \(F\) au schéma des zéros de \(\sigma_1,\ldots,\sigma_i\), alors \(h^0(F_i(-1))\le b_i\).}

\textbf{Proposition 1.6.} \emph{Soit \(F\) un \((b)\)-faisceau sur \(X\). Alors}~:
\begin{enumerate}[label=(\roman*)]
\item \emph{Pour toute suite ``assez générale'' \(\sigma\) de \(r\) sections \(\sigma_1,\ldots,\sigma_r\in H^0({O}_X(1))\), si \(F_{\sigma,i}\) désigne la restriction de \(F\) au schéma des zéros de \(\sigma_1,\ldots,\sigma_i\), alors \(h^0(F_{\sigma,i}(-1))\le b_i\), où \(0\le i\le r\).}
\item \emph{Tout sous-faisceau \(G\) de \(F\) est un \((b)\)-faisceau.}
\end{enumerate}

